% Chapter 1: Introduction
\section{Background and Motivation}

Computer science education faces significant challenges in the modern landscape. This section introduces the key motivations and context for this research.

\subsection{Educational Technology Landscape}

The rapid evolution of technology has transformed educational practices. Today's students require adaptive learning environments that respond to their individual needs and learning styles.

\subsection{Role of Ontologies}

Ontologies provide a structured, machine-readable representation of knowledge that can be leveraged to create intelligent educational systems. See Chapter~\ref{ch:ontology-design} for detailed discussion.

\vspace{1em}

\section{Research Problem and Objectives}

\subsection{Problem Statement}

Traditional computer science education often relies on one-size-fits-all pedagogical approaches. This research addresses the need for personalized, semantically-aware educational frameworks.

\subsection{Research Questions}

\begin{enumerate}
    \item How can ontologies be effectively used to model computer science education concepts?
    \item What system architecture best supports interactive, ontology-based learning?
    \item How do students respond to this interactive framework in practice?
\end{enumerate}

\vspace{1em}

\section{Contributions}

This dissertation makes the following key contributions:

\begin{enumerate}
    \item Development of a comprehensive CS education ontology
    \item Design of an interactive framework leveraging semantic technologies
    \item Empirical validation through user studies and performance metrics
\end{enumerate}

For more details on the framework design, refer to Part~\ref{part:framework-design-architecture} and the evaluation results in Chapter~\ref{ch:evaluation}.
